\section{Дополнительное задание}

В рамках дополнительного задания продемонстрирована работа с тегами,
командами отмены изменений revert и reset, перебазированием веток,
временным сохранением незавершённых изменений и серией патчей.

\subsection{Создание и применение тегов}

Для фиксации состояния репозитория, соответствующего версии 1.0, создан
аннотированный тег v1.0. Команде \verb|git tag -a| переданы идентификатор
коммита ba097ef и сообщение «Doing additional tasks». Наличие тега проверено
командами \verb|git tag| и \verb|git log --oneline --decorate|. Как видно
на рисунке~\ref{fig:additional-tag}, тег v1.0 указывает на выбранный коммит,
не изменяя историю ветки master.

\begin{figure}[H]
	\centering
	\reportScreenshot{figures/12-git-create-tag.png}
	\caption{Создание аннотированного тега v1.0}
	\label{fig:additional-tag}
\end{figure}

\subsection{Отмена изменений с помощью revert}

После создания ошибочного коммита «Probably this is a mistake...» на него
установлен тег v1.0.1.

Команда \verb|git revert v1.0.1| отменила это изменение. Тег использован
как ссылка на требуемый коммит. Git создал новый коммит, содержащий
обратные изменения, при этом
исходный ошибочный коммит остался в истории. Результат операции представлен
на рисунке~\ref{fig:additional-revert}.

\begin{figure}[H]
	\centering
	\reportScreenshot{figures/13-git-revert.png}
	\caption{Отмена коммита командой revert с применением тега}
	\label{fig:additional-revert}
\end{figure}

\subsection{Сброс состояния с помощью reset}

Для возврата ветки master к ранее отмеченной версии выполнена команда
\verb|git reset --hard v1.0|. В отличие от revert, команда reset переместила
указатель ветки непосредственно на коммит, отмеченный тегом v1.0, и привела
индекс и рабочий каталог к его состоянию. Коммиты, созданные после версии 1.0,
перестали входить в текущую историю ветки, что показано на
рисунке~\ref{fig:additional-reset}.

\begin{figure}[H]
	\centering
	\reportScreenshot{figures/14-git-reset.png}
	\caption{Жёсткий сброс ветки master до состояния, отмеченного тегом v1.0}
	\label{fig:additional-reset}
\end{figure}

\subsection{Перебазирование ветки (rebase)}

Для демонстрации перебазирования создана ветка readme, в которой добавлен
файл readme.md и выполнен коммит «Readme: added basic description».
В ветке master независимо создан файл с тем же именем, но другим содержимым,
после чего выполнено переключение обратно на ветку readme. Подготовка веток
к перебазированию показана на рисунке~\ref{fig:additional-rebase-prepare}.

\begin{figure}[H]
	\centering
	\reportScreenshot{figures/15-git-rebase.png}
	\caption{Подготовка веток master и readme к перебазированию}
	\label{fig:additional-rebase-prepare}
\end{figure}

При выполнении команды \verb|git rebase master| возник конфликт добавления
файла readme.md, так как файл был независимо создан в обеих ветках.
Конфликт разрешён вручную, исправленный файл добавлен в индекс командой
\verb|git add readme.md|, после чего перебазирование продолжено командой
\verb|git rebase --continue|. Итоговый граф истории на
рисунке~\ref{fig:additional-rebase-result} показывает, что коммит ветки
readme был пересоздан поверх актуального состояния master.

\begin{figure}[H]
	\centering
	\reportScreenshot{figures/16-git-rebase-conflict.png}
	\caption{Разрешение конфликта и завершение перебазирования ветки readme}
	\label{fig:additional-rebase-result}
\end{figure}

\subsection{Временное сохранение изменений (stash)}

В файл readme.md внесено незавершённое изменение, после чего рабочее состояние
сохранено командой \verb|git stash push -m| с сообщением
«WIP: новый крутой функционал».
Команда \verb|git stash list| подтвердила появление записи stash. После этого
в файл askname.h внесено срочное исправление и создан отдельный коммит
«hotfix». Команда \verb|git stash pop| восстановила отложенное изменение
readme.md и одновременно удалила использованную запись из stash. Затем
восстановленное изменение было зафиксировано отдельным коммитом. Полная
последовательность действий приведена на рисунке~\ref{fig:additional-stash}.

\begin{figure}[H]
	\centering
	\reportScreenshot{figures/17-git-stash.png}
	\caption{Сохранение изменений в stash и их последующее восстановление}
	\label{fig:additional-stash}
\end{figure}

\subsection{Создание и применение серии патчей}

В каталоге patches подготовлена серия из трёх патчей, соответствующих
последовательным изменениям «Feature part 1», «Feature part 2» и
«Feature part 3». Каждый патч применён командой \verb|git am|. Git сохранил
авторство и сообщения исходных коммитов; выведенные предупреждения относятся
к завершающим пробелам в строках патчей и не препятствуют их применению.
Процесс показан на рисунке~\ref{fig:additional-patches-apply}.

\begin{figure}[H]
	\centering
	\reportScreenshot{figures/18-git-patches.png}
	\caption{Последовательное применение серии патчей командой git am}
	\label{fig:additional-patches-apply}
\end{figure}

После применения патчей выполнена команда \verb|git log --oneline|.
В истории ветки master появились три последовательных коммита Feature part 1,
Feature part 2 и Feature part 3, что подтверждает корректное применение всей
серии. Итоговая история показана на
рисунке~\ref{fig:additional-patches-result}.

\begin{figure}[H]
	\centering
	\reportScreenshot{figures/19-git-log.png}
	\caption{История репозитория после применения серии патчей}
	\label{fig:additional-patches-result}
\end{figure}

Для создания серии использована команда \verb|git format-patch|
с параметрами \verb|-6 -o patches|. В каталоге patches сформировано
шесть файлов, по одному для каждого из шести последних коммитов. Имена файлов
содержат порядковый номер и сообщение соответствующего коммита. Результат
создания серии представлен на рисунке~\ref{fig:additional-patches-create}.

\begin{figure}[H]
	\centering
	\reportScreenshot{figures/20-git-patches-list.png}
	\caption{Создание серии патчей командой git format-patch}
	\label{fig:additional-patches-create}
\end{figure}

Таким образом, в дополнительном задании продемонстрированы два способа
отмены изменений, использование тегов как именованных указателей на коммиты,
перенос коммитов между основаниями с помощью rebase, временное сохранение
рабочего состояния в stash, а также перенос серии изменений посредством
патчей.
